% GENERATED FILE. Edit canonical lessons, then run npm run generate:books.
% canonical-source-hash: fnv1a64:49135e73379709d8
% canonical-lessons: TA-C490-nimir, TA-C490-kuni, TA-C490-cay, TA-C490-utaru, TA-C490-no

\chapter{Bending and Leaning}
\label{ch:ta-bending-and-leaning}
\hlchaptermodality{490}

\begin{chapteropening}
\textbf{By the end of this chapter:} \emph{I can talk about straightening up, bending, leaning, shaking off and aching.}

Complete the last lesson of chapter 490: I can talk about straightening up, bending, leaning, shaking off and aching.
\end{chapteropening}

\section[\texorpdfstring{\ta{நிமிர்} (nimir)}{நிமிர் (nimir)}]{\ta{நிமிர்} (nimir) — to straighten up}

\label{lesson:TA-C490-nimir}



\begin{quote}
Before the new one: say the Tamil for to wake up, then the Tamil for to perspire.
\end{quote}

\subsection*{You'll want to know: \ta{நிமிர்}}
\textbf{\ta{நிமிர்}} — \emph{nimir} — \textquotedblleft{}to straighten up\textquotedblright{}.

\begin{quote}
\ta{வேலை} \ta{முடிந்ததும்} \ta{அவள்} \ta{நிமிர்ந்து} \ta{உட்கார்ந்தாள்}.
\end{quote}

\begin{quote}
\emph{vēlai muṭintatum avaḷ nimirntu uṭkārntāḷ.}
\end{quote}

\begin{quote}
When the work was done she sat up straight.
\end{quote}

Past (I …): \ta{நிமிர்ந்தேன்}, \emph{nimirntēṉ}.

\ta{நிமிர்ந்து} \ta{நில்}, to stand tall.

\begin{cousinweb}
\textbf{\ta{நிமிர்}} (\emph{nimir}) is native Dravidian (DEDR 2922), to stand or sit upright.
\end{cousinweb}

\begin{grammarlens}[title={What you've built}]
One more word for describing what you lived and felt.
\end{grammarlens}

\subsection*{Your turn}
\begin{itemize}
\raggedright
  \item \emph{Say it:} \emph{nimir}
  \item \emph{Say it:} \emph{nimir}, once more
  \item \emph{Say it:} \emph{vēlai muṭintatum avaḷ nimirntu uṭkārntāḷ.} — When the work was done she sat up straight.
  \item \emph{Recall:} say the Tamil for to wake up, then the Tamil for to perspire, then say \emph{nimir} again
\end{itemize}

\begin{tcolorbox}[breakable,colback=teal!4,colframe=teal!35!black,title={Before you move on}]
What does \ta{நிமிர்} mean? (\textquotedblleft{}To straighten up\textquotedblright{}.) Say it once more.
\end{tcolorbox}
\section[\texorpdfstring{\ta{குனி} (kuṉi)}{குனி (kuṉi)}]{\ta{குனி} (kuṉi) — to bend down; to stoop}

\label{lesson:TA-C490-kuni}



\begin{quote}
Before the new one: say the Tamil for to perspire, then the Tamil for to straighten up.
\end{quote}

\subsection*{You'll want to know: \ta{குனி}}
\textbf{\ta{குனி}} — \emph{kuṉi} — \textquotedblleft{}to bend down; to stoop\textquotedblright{}.

\begin{quote}
\ta{கீழே} \ta{விழுந்த} \ta{பேனாவை} \ta{எடுக்க} \ta{நான்} \ta{குனிந்தேன்}.
\end{quote}

\begin{quote}
\emph{kīḻē viḻunta pēṉāvai eṭukka nāṉ kuṉintēṉ.}
\end{quote}

\begin{quote}
I bent down to pick up the pen that had fallen.
\end{quote}

Past (I …): \ta{குனிந்தேன்}, \emph{kuṉintēṉ}.

Its opposite is \ta{நிமிர்}.

\begin{cousinweb}
\textbf{\ta{குனி}} (\emph{kuṉi}) is native Dravidian (DEDR 1927), in the entry of \ta{கூன்} \textquotedblleft{}a bend, a hump\textquotedblright{}.
\end{cousinweb}

\begin{grammarlens}[title={What you've built}]
One more word for describing what you lived and felt.
\end{grammarlens}

\subsection*{Your turn}
\begin{itemize}
\raggedright
  \item \emph{Say it:} \emph{kuṉi}
  \item \emph{Say it:} \emph{kuṉi}, once more
  \item \emph{Say it:} \emph{kīḻē viḻunta pēṉāvai eṭukka nāṉ kuṉintēṉ.} — I bent down to pick up the pen that had fallen.
  \item \emph{Recall:} say the Tamil for to perspire, then the Tamil for to straighten up, then say \emph{kuṉi} again
\end{itemize}

\begin{tcolorbox}[breakable,colback=teal!4,colframe=teal!35!black,title={Before you move on}]
What does \ta{குனி} mean? (\textquotedblleft{}To bend down; to stoop\textquotedblright{}.) Say it once more.
\end{tcolorbox}
\section[\texorpdfstring{\ta{சாய்} (cāy)}{சாய் (cāy)}]{\ta{சாய்} (cāy) — to lean; to tilt}

\label{lesson:TA-C490-cay}



\begin{quote}
Before the new one: say the Tamil for to straighten up, then the Tamil for to bend down.
\end{quote}

\subsection*{You'll want to know: \ta{சாய்}}
\textbf{\ta{சாய்}} — \emph{cāy} — \textquotedblleft{}to lean; to tilt\textquotedblright{}.

\begin{quote}
\ta{களைப்பாக} \ta{இருந்ததால்} \ta{அவள்} \ta{சுவரில்} \ta{சாய்ந்தாள்}.
\end{quote}

\begin{quote}
\emph{kaḷaippāka iruntatāl avaḷ cuvaril cāyntāḷ.}
\end{quote}

\begin{quote}
She was tired, so she leaned against the wall.
\end{quote}

Past (I …): \ta{சாய்ந்தேன்}, \emph{cāyntēṉ}.

What you lean on takes -\ta{இல்}: \ta{சுவரில்} \ta{சாய்}.

\begin{cousinweb}
\textbf{\ta{சாய்}} (\emph{cāy}) is native Dravidian (DEDR 2456). DEDR has other \ta{சாய்} entries, such as \textquotedblleft{}colour\textquotedblright{}, that are different words.
\end{cousinweb}

\begin{grammarlens}[title={What you've built}]
One more word for describing what you lived and felt.
\end{grammarlens}

\subsection*{Your turn}
\begin{itemize}
\raggedright
  \item \emph{Say it:} \emph{cāy}
  \item \emph{Say it:} \emph{cāy}, once more
  \item \emph{Say it:} \emph{kaḷaippāka iruntatāl avaḷ cuvaril cāyntāḷ.} — She was tired, so she leaned against the wall.
  \item \emph{Recall:} say the Tamil for to straighten up, then the Tamil for to bend down, then say \emph{cāy} again
\end{itemize}

\begin{tcolorbox}[breakable,colback=teal!4,colframe=teal!35!black,title={Before you move on}]
What does \ta{சாய்} mean? (\textquotedblleft{}To lean; to tilt\textquotedblright{}.) Say it once more.
\end{tcolorbox}
\section[\texorpdfstring{\ta{உதறு} (utaṟu)}{உதறு (utaṟu)}]{\ta{உதறு} (utaṟu) — to shake off; to shake out}

\label{lesson:TA-C490-utaru}



\begin{quote}
Before the new one: say the Tamil for to bend down, then the Tamil for to lean.
\end{quote}

\subsection*{You'll want to know: \ta{உதறு}}
\textbf{\ta{உதறு}} — \emph{utaṟu} — \textquotedblleft{}to shake off; to shake out\textquotedblright{}.

\begin{quote}
\ta{அவள்} \ta{போர்வையை} \ta{உதறி} \ta{மடித்தாள்}.
\end{quote}

\begin{quote}
\emph{avaḷ pōrvaiyai utaṟi maṭittāḷ.}
\end{quote}

\begin{quote}
She shook out the blanket and folded it.
\end{quote}

Past (I …): \ta{உதறினேன்}, \emph{utaṟiṉēṉ}.

Also to brush aside: \ta{பயத்தை} \ta{உதறினான்}, he shook off his fear.

\begin{cousinweb}
\textbf{\ta{உதறு}} (\emph{utaṟu}) is native Dravidian (DEDR 613), of shaking out a cloth or casting off a tie.
\end{cousinweb}

\begin{grammarlens}[title={What you've built}]
One more word for describing what you lived and felt.
\end{grammarlens}

\subsection*{Your turn}
\begin{itemize}
\raggedright
  \item \emph{Say it:} \emph{utaṟu}
  \item \emph{Say it:} \emph{utaṟu}, once more
  \item \emph{Say it:} \emph{avaḷ pōrvaiyai utaṟi maṭittāḷ.} — She shook out the blanket and folded it.
  \item \emph{Recall:} say the Tamil for to bend down, then the Tamil for to lean, then say \emph{utaṟu} again
\end{itemize}

\begin{tcolorbox}[breakable,colback=teal!4,colframe=teal!35!black,title={Before you move on}]
What does \ta{உதறு} mean? (\textquotedblleft{}To shake off; to shake out\textquotedblright{}.) Say it once more.
\end{tcolorbox}
\section[\texorpdfstring{\ta{நோ} (nō)}{நோ (nō)}]{\ta{நோ} (nō) — to ache; to hurt}

\label{lesson:TA-C490-no}



\begin{quote}
Before the new one: say the Tamil for to lean, then the Tamil for to shake off.
\end{quote}

\subsection*{You'll want to know: \ta{நோ}}
\textbf{\ta{நோ}} — \emph{nō} — \textquotedblleft{}to ache; to hurt\textquotedblright{}.

\begin{quote}
\ta{நீண்ட} \ta{நேரம்} \ta{நின்றதால்} \ta{கால்} \ta{நொந்தது}.
\end{quote}

\begin{quote}
\emph{nīṇṭa nēram niṉṟatāl kāl nontatu.}
\end{quote}

\begin{quote}
My legs ached from standing a long time.
\end{quote}

Past (it …): \ta{நொந்தது}, \emph{nontatu}.

More literary than \ta{வலி}, the everyday verb. Also of the heart: \ta{மனம்} \ta{நொந்தது}, I was hurt. The past shortens the vowel: \ta{நொந்தது}.

\begin{cousinweb}
\textbf{\ta{நோ}} (\emph{nō}) is native Dravidian (DEDR 3793); \ta{நோய்} \textquotedblleft{}disease\textquotedblright{} comes from the same root.
\end{cousinweb}

\begin{grammarlens}[title={What you've built}]
One more word for describing what you lived and felt.
\end{grammarlens}

\subsection*{Your turn}
\begin{itemize}
\raggedright
  \item \emph{Say it:} \emph{nō}
  \item \emph{Say it:} \emph{nō}, once more
  \item \emph{Say it:} \emph{nīṇṭa nēram niṉṟatāl kāl nontatu.} — My legs ached from standing a long time.
  \item \emph{Recall:} say the Tamil for to lean, then the Tamil for to shake off, then say \emph{nō} again
\end{itemize}

\begin{tcolorbox}[breakable,colback=teal!4,colframe=teal!35!black,title={Before you move on}]
What does \ta{நோ} mean? (\textquotedblleft{}To ache; to hurt\textquotedblright{}.) Say it once more.
\end{tcolorbox}
